\documentclass[12pt]{article}

\usepackage{fullpage}
\usepackage{graphicx}
\usepackage{amssymb}
\usepackage{amsmath}
\usepackage[none]{hyphenat}
\usepackage{parskip}
\usepackage[spanish]{babel}
\usepackage[utf8]{inputenc}
\usepackage{hyperref}
\usepackage{fancyhdr}
\usepackage{tasks}
\usepackage{mdframed}
\usepackage{xcolor}
\usepackage{pgfplots}
\usepackage[makeroom]{cancel}
\usepackage{multicol}
\usepackage[shortlabels]{enumitem}
\usepackage{stackrel}
\usepackage{tkz-tab}
\usepackage{xpatch}
\usepackage{tkz-euclide}
\usetkzobj{all}
\xpatchcmd{\tkzTabLine}{$0$}{$\bullet$}{}{}

\setlength{\headheight}{10pt}
\setlength{\headsep}{10pt}
\pagestyle{fancy}
\rhead{\ayudantia \ - \alumno}
\tikzset{t style/.style={style=solid}}

\newcommand*{\mybox}[2]{\colorbox{#1!30}{\parbox{.98\linewidth}{#2}}}

\newenvironment{solucion}
{\begin{mdframed}[backgroundcolor=black!10]
		{\bf Solución:}\\
	}
	{
	\end{mdframed}
}

\newenvironment{alternativas}[1]
{\begin{multicols}{#1}
		\begin{enumerate}[a)]
		}
		{
		\end{enumerate}
	\end{multicols}
}

\newenvironment{preguntas}
{\begin{enumerate}\itemsep12pt
	}
	{
	\end{enumerate}
}

\newcommand{\ayudantia}{{\sc Ayudantía 5}}
\newcommand{\tituloayu}{Series III}
\newcommand{\fecha}{10 de septiembre de 2019}
\newcommand{\sigla}{MAT1620}
\newcommand{\nombre}{Cálculo II}
\newcommand{\profesor}{Wolfgang Rivera}
\newcommand{\ano}{2019}
\newcommand{\semestre}{2}
\newcommand{\mail}{mat1620@ifcastaneda.cl}
\newcommand{\alumno}{Ignacio Castañeda - \mail}

\newcommand{\ev}{\Big|}
\newcommand{\ra}{\rightarrow}
\newcommand{\lra}{\leftrightarrow}
\newcommand{\N}{\mathbb{N}}
\newcommand{\R}{\mathbb{R}}
\newcommand{\Exp}[1]{\mathcal{E}_{#1}}
\newcommand{\List}[1]{\mathcal{L}_{#1}}
\newcommand{\EN}{\Exp{\N}}
\newcommand{\LN}{\List{\N}}
\newcommand{\comment}[1]{}
\newcommand{\lb}{\\~\\}
\newcommand{\eop}{_{\square}}
\newcommand{\hsig}{\hat{\sigma}}
\newcommand{\widesim}[2][1.5]{
	\mathrel{\overset{#2}{\scalebox{#1}[1]{$\sim$}}}
}
\newcommand{\wsim}{\widesim{}}
\newcommand{\lh}{\stackrel{L'H}{=}}

\begin{document}
\thispagestyle{empty}

\begin{minipage}{2cm}
	\includegraphics[width=2cm]{../../../../img/logo.pdf}
	\vspace{0.5cm}
\end{minipage}
\begin{minipage}{\linewidth}
	\begin{tabular}{lrl}
		{\scriptsize\sc Pontificia Universidad Catolica de Chile} & \hspace*{0.7in}Curso: &
		\sigla  - \nombre\\
		{\sc Facultad de Matemáticas}&
		Profesor: & \profesor \\
		{\sc Semestre \ano-\semestre} & Ayudante: & {Ignacio Castañeda}\\
		& {Mail:} & \texttt{\mail}
	\end{tabular}
\end{minipage}

\vspace{-10mm}
\begin{center}
	{\LARGE\bf \ayudantia}\\
	\vspace{0.1cm}
	{\tituloayu}\\
	\vspace{0.1cm}
	\fecha\\
	\vspace{0.4cm}
\end{center}

\begin{preguntas}
\item Escribir las siguientes funciones como una serie de potencias
\begin{tasks}(2)
\task $f(x) = \dfrac{1}{x^2-4x+20}$
\task $f(x) = ln(1+x)$
\end{tasks}
\begin{solucion}
Recordemos que $$\dfrac{1}{1-x} = \sum\limits_0^{\infty} x^n$$
\begin{enumerate}[a)]
\item $f(x) = \dfrac{1}{x^2-4x+20}$\\
\\
Debemos reescribir la función para que se asemeje a la formula de más arriba. Para esto, debemos completar cuadrados, esto es
$$f(x) = \dfrac{1}{x^2-4x+20}
= \dfrac{1}{16 + (x^2-4x+4)}
= \dfrac{1}{16 + (x-2)^2}$$
Ahora, debemos convertir el 16 en un 1, para esto, basta con factorizar, es decir
$$ f(x)
= \dfrac{1}{16 + (x-2)^2}
= \dfrac{1}{16}\cdot \dfrac{1}{1 + \frac{(x-2)^2}{16}}
= \dfrac{1}{16} \cdot\dfrac{1}{1 + \left(\frac{x-2}{4}\right)^2}
$$
Por último, necesitamos un signo menos luego del 1.
$$ f(x)
= \dfrac{1}{16} \cdot\dfrac{1}{1 + \left(\frac{x-2}{4}\right)}
= \dfrac{1}{16} \cdot\dfrac{1}{1 + \left(-\left(\frac{x-2}{4}\right)^2\right)}
$$
Aplicando la formula,
$$ f(x)
= \dfrac{1}{16} \cdot\dfrac{1}{1 + \left(-\left(\frac{x-2}{4}\right)^2\right)}
= \sum\limits_{n=0}^{\infty} \left(-\left(\frac{x-2}{4}\right)^2\right)^n
= \sum\limits_{n=0}^{\infty} (-1)^n \dfrac{(x-2)^{2n}}{16^n}
$$
\item $f(x) = ln(1+x)$\\
\\
Es evidente que no podemos convertir la función directamente a una serie, por lo que probaremos derivandola antes.
		$$f'(x) = \dfrac{1}{1+x}$$
		Esta función si podemos convertirla en una serie usando $\dfrac{1}{1-x} = \sum\limits_{n=0}^{\infty} x^n$. Entonces,
		$$f'(x) = \dfrac{1}{1-(-x)} 
		= \sum\limits_{n=0}^{\infty} (-x)^n
		= \sum\limits_{n=0}^{\infty} (-1)^n x^n$$
		Lo único que nos queda por hacer es integrar esto. Recordemos que para integrar una serie de potencias, basta con integrar su termino general.
		
		Finalmente, 
		$$f(x) 
		= \int f'(x) dx
		= \int \sum\limits_{n=0}^{\infty} (-1)^n x^n dx
		= \sum\limits_{n=0}^{\infty} \int (-1)^n x^n dx
		= \sum\limits_{n=0}^{\infty} (-1)^n \dfrac{x^{n+1}}{n+1}$$
\end{enumerate}
\end{solucion}
\item Expresar las siguientes series de potencias como una función
\begin{tasks}(2)
\task $\sum\limits_{n=2}^\infty \dfrac{x^n}{2^{n-1}}$
\task $\sum\limits_{n=1}^\infty \dfrac{(-1)^n (3x+1)^{n-1}}{5^n}$
\task $\sum\limits_{n=0}^\infty \dfrac{x^{2n}}{4^n}$
\task $\sum\limits_{n=1}^\infty \dfrac{(2x-3)^n}{n2^n}$
\end{tasks}
\begin{solucion}

	Para hacer esto debemos seguir los siguientes pasos
	\begin{itemize}
		\item Hacer que la serie comience del 0
		\item Mover todo hacia afuera de tal forma de que todos los exponentes sean $n$
		\item Aplicar $\dfrac{1}{1-x} = \sum\limits_0^{\infty} x^n$
		\item Simplificar
	\end{itemize}
\begin{enumerate}[a)]
\item $\sum\limits_{n=2}^\infty \dfrac{x^n}{2^{n-1}}
= \sum\limits_{n=0}^\infty \dfrac{x^{n+2}}{2^{n+1}}
= \dfrac{x^2}{2} \sum\limits_{n=0}^\infty \left(\dfrac{x}{2}\right)^n
= \dfrac{x^2}{2} \dfrac{1}{1-\frac{x}{2}}
= \dfrac{x^2}{2-x}
$
\item $\sum\limits_{n=1}^\infty \dfrac{(-1)^n (3x+1)^{n-1}}{5^n}
= \sum\limits_{n=0}^\infty \dfrac{(-1)^{n+1} (3x+1)^{n}}{5^{n+1}}
= -\dfrac{1}{5} \sum\limits_{n=0}^\infty \dfrac{(-1)^{n} (3x+1)^{n}}{5^{n}}$\\
$
= -\dfrac{1}{5} \sum\limits_{n=0}^\infty \left(\dfrac{-(3x+1)}{5}\right)^n
= -\dfrac{1}{5} \cdot \dfrac{1}{1-\left(\frac{-(3x+1)}{5}\right)}
= \dfrac{-1}{3x+6}
$
\item $\sum\limits_{n=0}^\infty \dfrac{x^{2n}}{4^n}
= \sum\limits_{n=0}^\infty \left(\dfrac{x^2}{4}\right)^n
= \dfrac{1}{1-\frac{x}{4}}
= \dfrac{4}{4-x^2}	
$
\item Sea 
		$$f(x) = \sum\limits_{n=1}^\infty \dfrac{(2x-3)^n}{n2^n}$$
		En este caso, lo que nos molesta es el $n$ del denominador, por lo que intentaremos derivando la serie. De manera análoga, para hacer esto solo tenemos que derivar el término general de esta, esto es,
		$$f'(x) 
		= \sum\limits_{n=1}^\infty \dfrac{n(2x-3)^{n-1}\cdot 2}{n2^n}
		= \sum\limits_{n=1}^\infty \dfrac{(2x-3)^{n-1}}{2^{n-1}}
		= \sum\limits_{n=1}^\infty \left(x-\dfrac{3}{2}\right)^{n-1}
		= \sum\limits_{n=0}^\infty \left(x-\dfrac{3}{2}\right)^{n}$$
		Ahora que no tenemos problema,
		$$f'(x) 
		= \dfrac{1}{1-\left(x-\dfrac{3}{2}\right)}
		= \dfrac{1}{1-x+\dfrac{3}{2}}
		= \dfrac{1}{\dfrac{5}{2}-x}
		= \dfrac{1}{\dfrac{5-2x}{2}}
		= \dfrac{2}{5-2x}$$
		Integrando,
		$$f(x) 
		= \int \dfrac{2}{5-2x} dx
		= -ln (5-2x) + c$$
		Para encontrar la constante, debemos evaluar la función en un punto donde conozcamos el valor de la serie que representa. La opción trivial es evaluarla donde el termino general se hacer cero, esto sería en $x = \dfrac{3}{2}$. Aquí,
		$$f\left(\dfrac{3}{2}\right) = \sum\limits_{n=1}^\infty 0 = 0$$
		Luego,
		$$f\left(\dfrac{3}{2}\right) = -ln(5-3) + c = 0 \ra c = ln(2)$$
		Finalmente,
		$$f(x) = -ln(5-2x) + ln(2) = ln\left(\dfrac{2}{5-2x}\right)$$
\end{enumerate}
\end{solucion}
\item Determinar, utilizando series de potencias, el valor de
	$$\sum\limits_{n=1}^\infty \frac{1}{3^n}$$
\begin{solucion}
Recordemos que $$\dfrac{1}{1-x} = \sum\limits_0^{\infty} x^n$$
		Luego, debemos buscar una serie de potencias que se asemeje a lo que estamos buscando, esta sería
		$$f(x) = \sum\limits_{n=1}^\infty x^{n}$$
		Entonces, debemos encontrar $f\left(\dfrac{1}{3}\right)$
		
		Tenemos que
		$$f(x) 
		= \sum\limits_{n=1}^\infty x^{n} 
		= \sum\limits_{n=0}^\infty x^{n+1}
		= x\sum\limits_{n=0}^\infty x^{n}
		= x \dfrac{1}{1-x}
		= \dfrac{x}{1-x}$$
		Como $x = \dfrac{1}{3}$ esta dentro del radio de convergencia de la serie, podremos evaluar la seríe ahí. Finalmente,
		$$\sum\limits_{n=1}^\infty \frac{1}{3^n} = f\left(\dfrac{1}{3}\right) = \dfrac{\dfrac{1}{3}}{1-\dfrac{1}{3}}
		= \dfrac{\dfrac{1}{3}}{\dfrac{2}{3}} = \dfrac{1}{2}$$
\end{solucion}
\item Determinar el valor de la siguiente serie
	$$ \sum\limits_{n=2}^{\infty}\dfrac{n^2+n}{3^{n-1}} $$
\begin{solucion}
En primer lugar, notemos que la serie se puede escribir como
$$ \sum\limits_{n=2}^{\infty}\dfrac{n(n+1)}{3^{n-1}} $$
Lo que nos da problema es el $n(n+1)$, sin embargo, notemos que este esta en el numerador, por lo que debemos convertir la serie en una serie de potencias apropiada para poder integrarla.	\\
\\
Sea
$$f(x) = \sum\limits_{n=2}^{\infty}\dfrac{n(n+1)}{3^{n-1}} x^{n-1} $$
Debemos encontrar el valor de $f(1)$\\

Integramos,
$$\int f(x)dx 
= \sum\limits_{n=2}^{\infty}\dfrac{\cancel{n}(n+1)}{3^{n-1}} \dfrac{x^n}{\cancel{n}} 
= \sum\limits_{n=2}^{\infty}\dfrac{n+1}{3^{n-1}} x^n$$
Integramos  nuevamente,
$$\iint f(x)dxdx 
= \sum\limits_{n=2}^{\infty}\dfrac{\cancel{n+1}}{3^{n-1}} \dfrac{x^{n+1}}{\cancel{n+1}} 
= \sum\limits_{n=2}^{\infty}\dfrac{x^{n+1}}{3^{n-1}}$$
Notemos que
$$\sum\limits_{n=2}^{\infty}\dfrac{x^{n+1}}{3^{n-1}}
= \sum\limits_{n=0}^{\infty}\dfrac{x^{n+3}}{3^{n+1}}
= \dfrac{x^3}{3}\sum\limits_{n=0}^{\infty}\left(\dfrac{x}{3}\right)^n
= \dfrac{x^3}{3}\dfrac{1}{1-\dfrac{x}{3}}
= \dfrac{x^3}{3-x}$$
Por lo que
$$\iint f(x)dxdx 
= \dfrac{x^3}{3-x}$$
Derivamos,
$$\int f(x)dx
= \left(\dfrac{x^3}{3-x}\right)' 
= \dfrac{9x^2-2x^3}{(3-x)^2}
$$
Derivamos nuevamente,
$$f(x)
= \left(\dfrac{9x^2-2x^3}{(3-x)^2}\right)' 
= \dfrac{(18x-6x^2)(3-x)^2 + 2(3-x)(9x^2-2x^3)}{(3-x)^4}
$$
Finalmente,
$$f(1)
= \dfrac{(18-6)(3-1)^2 + 2(3-1)(9-2)}{(3-1)^4}
=\dfrac{19}{4}
$$
Por lo que
$$ \sum\limits_{n=2}^{\infty}\dfrac{n(n+1)}{3^{n-1}} 
=\dfrac{19}{4}$$
\end{solucion}
\item Para cada función, encontrar la serie de Maclaurin que la representa
\begin{tasks}(2)
\task $f(x) = cos(x)$
\task $f(x) = \dfrac{1}{1-x}$
\end{tasks}
\begin{solucion}

\begin{enumerate}[a)]
\item 
\item 
\end{enumerate}
\end{solucion}
\item Encontrar una aproximación de $\sqrt[]{101}$ con un error máximo de $10^{-3}$
\begin{solucion}

\end{solucion}
\end{preguntas}
\end{document}